\documentclass[12pt,a4paper]{ctexart}
\usepackage{amsmath,amssymb}
\usepackage{booktabs}
\usepackage{geometry}
\geometry{left=2.5cm,right=2.5cm,top=2.5cm,bottom=2.5cm}
\usepackage{array}

\newcommand{\Var}{\operatorname{Var}}
\newcommand{\Cov}{\operatorname{Cov}}
\newcommand{\VIF}{\operatorname{VIF}}

% 只对 section 进行自动编号，subsection 及以下不编号
\setcounter{secnumdepth}{1}

\pagestyle{plain}

\title{\textbf{多重共线性产生的后果——极详细推导与讲解}}
\author{}
\date{}

\begin{document}

\maketitle

\section*{引言}
\addcontentsline{toc}{section}{引言}

本文以初学者能跟上的细致程度，推导和讲解“多重共线性”对回归分析造成的全部后果。内容分为两大部分：\textbf{完全多重共线性}的后果和\textbf{不完全多重共线性}的后果。每一步代数变形、每一个直觉理解、每一个公式背后的“为什么”都给出解释，绝不存在跳步。

% 让接下来的 section 从 0 开始编号
\setcounter{section}{0}
\section{超级基础的准备知识}

\subsection{0.1 我们在研究什么问题？}

在多元线性回归中，我们用多个解释变量（X）来解释一个被解释变量（Y）。比如，用“工业增加值”和“居民消费”两个变量来解释“财政收入”。

\textbf{一个自然的问题}：如果两个解释变量之间本身就存在很强的线性关系（比如工业增加值和居民消费几乎同步增长），会出什么问题？

这就是“多重共线性”要讨论的核心。

\subsection{0.2 什么是多重共线性？}

\begin{itemize}
    \item \textbf{直观说法}：解释变量之间“彼此长得太像”，一个变量几乎可以用另一个变量乘上一个常数来表示。
    \item \textbf{严格说法}：如果某个解释变量可以由其他解释变量\textbf{精确地}线性表示，就叫做“完全多重共线性”；如果只是\textbf{近似地}线性表示，就叫做“不完全多重共线性”。
\end{itemize}

\subsection{0.3 本文使用的模型}

为了具体说明问题，我们始终以\textbf{两个解释变量的回归模型}为例：

\begin{equation}
Y_i = \beta_1 + \beta_2 X_{2i} + \beta_3 X_{3i} + u_i \tag{0.1}
\end{equation}

其中：
\begin{itemize}
    \item $Y_i$：第 $i$ 个观测的被解释变量。
    \item $X_{2i}$：第 $i$ 个观测的第一个解释变量。
    \item $X_{3i}$：第 $i$ 个观测的第二个解释变量。
    \item $\beta_1$：截距项（常数项）。
    \item $\beta_2$：$X_2$ 对 $Y$ 的偏回归系数——含义是“在 $X_3$ 不变的条件下，$X_2$ 每变动1单位，$Y$ 平均变动 $\beta_2$ 单位”。
    \item $\beta_3$：$X_3$ 对 $Y$ 的偏回归系数——含义是“在 $X_2$ 不变的条件下，$X_3$ 每变动1单位，$Y$ 平均变动 $\beta_3$ 单位”。
    \item $u_i$：随机误差项。
\end{itemize}

\subsection{0.4 离差形式}

为了推导方便，我们把模型写成离差形式。令：
\begin{itemize}
    \item $y_i = Y_i - \bar{Y}$（$Y$ 的离差）
    \item $x_{2i} = X_{2i} - \bar{X}_2$（$X_2$ 的离差）
    \item $x_{3i} = X_{3i} - \bar{X}_3$（$X_3$ 的离差）
\end{itemize}

离差形式的回归模型为：

\begin{equation}
y_i = \beta_2 x_{2i} + \beta_3 x_{3i} + (u_i - \bar{u}) \tag{4.4}
\end{equation}

\textbf{暂停理解一下}：离差形式的好处是消去了截距项 $\beta_1$，使得公式更简洁。原来的模型有3个参数（$\beta_1, \beta_2, \beta_3$），离差形式只有2个（$\beta_2, \beta_3$），计算更方便。

\subsection{0.5 离差的重要性质}

\begin{equation}
\sum_{i=1}^{n} x_{2i} = \sum_{i=1}^{n} (X_{2i} - \bar{X}_2) = \sum X_{2i} - n\bar{X}_2 = n\bar{X}_2 - n\bar{X}_2 = 0
\end{equation}

同理 $\sum x_{3i} = 0$，$\sum y_i = 0$。这个性质后面会反复用到。

\subsection{0.6 OLS估计量的方差公式来源}

在多元回归中，参数估计量的方差-协方差矩阵为：

\begin{equation}
\text{Var-Cov}(\hat{\boldsymbol{\beta}}) = \sigma^2 (\mathbf{X}'\mathbf{X})^{-1}
\end{equation}

这个公式告诉我们：参数估计量的精度取决于 $(\mathbf{X}'\mathbf{X})^{-1}$ 这个矩阵。如果这个逆矩阵出了问题（不存在、或者元素很大），估计量就出问题了。\textbf{多重共线性的所有后果，本质上都源于这个逆矩阵出了问题。}

\section{第一部分：完全多重共线性产生的后果}

\subsection{后果1：参数的估计值不确定}

\textbf{目标}：证明当 $X_2$ 和 $X_3$ 完全共线性时（即 $X_{2i} = \lambda X_{3i}$），OLS估计量 $\hat{\beta}_2$ 和 $\hat{\beta}_3$ 变成 0/0 的不定式，无法确定。

\subsubsection{第一步：回顾OLS估计公式}

由第三章的结果，对离差形式的模型 $y_i = \beta_2 x_{2i} + \beta_3 x_{3i} + (u_i - \bar{u})$，其OLS估计量为：

\begin{equation}
\hat{\beta}_2 = \frac{(\sum y_i x_{2i})(\sum x_{3i}^2) - (\sum y_i x_{3i})(\sum x_{2i} x_{3i})}{(\sum x_{2i}^2)(\sum x_{3i}^2) - (\sum x_{2i} x_{3i})^2} \tag{4.5}
\end{equation}

\begin{equation}
\hat{\beta}_3 = \frac{(\sum y_i x_{3i})(\sum x_{2i}^2) - (\sum y_i x_{2i})(\sum x_{2i} x_{3i})}{(\sum x_{2i}^2)(\sum x_{3i}^2) - (\sum x_{2i} x_{3i})^2} \tag{4.6}
\end{equation}

\textbf{暂停理解一下}：这两个公式的分母是完全一样的！都是 $(\sum x_{2i}^2)(\sum x_{3i}^2) - (\sum x_{2i} x_{3i})^2$。记住这一点，后面会反复用到。

\subsubsection{第二步：引入完全共线性条件}

现在假设 $X_2$ 和 $X_3$ 完全共线性，即存在一个非零常数 $\lambda$，使得：

\[
X_{2i} = \lambda X_{3i} \quad \text{对所有 } i
\]

这意味着什么呢？$X_2$ 不过是 $X_3$ 的一个固定倍数。比如，如果 $X_2$ 是“以厘米衡量的身高”，$X_3$ 是“以米衡量的身高”，那么 $X_2 = 100 X_3$（$\lambda = 100$）。它们本质上是同一个信息！

既然 $X_{2i} = \lambda X_{3i}$，那么离差之间也满足同样的关系：

\[
x_{2i} = X_{2i} - \bar{X}_2 = \lambda X_{3i} - \lambda \bar{X}_3 = \lambda(X_{3i} - \bar{X}_3) = \lambda x_{3i}
\]

所以：\textbf{$x_{2i} = \lambda x_{3i}$}。

\subsubsection{第三步：将 $x_{2i} = \lambda x_{3i}$ 代入分子和分母}

\textbf{先处理分母}（两个公式共用）：

分母 $= (\sum x_{2i}^2)(\sum x_{3i}^2) - (\sum x_{2i} x_{3i})^2$

把 $x_{2i} = \lambda x_{3i}$ 代入：
\begin{itemize}
    \item $\sum x_{2i}^2 = \sum (\lambda x_{3i})^2 = \lambda^2 \sum x_{3i}^2$
    \item $\sum x_{2i} x_{3i} = \sum (\lambda x_{3i})(x_{3i}) = \lambda \sum x_{3i}^2$
\end{itemize}

所以分母变成：

\[
\text{分母} = (\lambda^2 \sum x_{3i}^2)(\sum x_{3i}^2) - (\lambda \sum x_{3i}^2)^2
\]

展开：

\[
= \lambda^2 (\sum x_{3i}^2)(\sum x_{3i}^2) - \lambda^2 (\sum x_{3i}^2)^2
\]
\[
= \lambda^2 (\sum x_{3i}^2)^2 - \lambda^2 (\sum x_{3i}^2)^2 = 0
\]

\textbf{分母等于0！} 这意味着无论分子是什么，整个分数都没有意义。

\textbf{再处理 $\hat{\beta}_2$ 的分子}：

分子 $= (\sum y_i x_{2i})(\sum x_{3i}^2) - (\sum y_i x_{3i})(\sum x_{2i} x_{3i})$

把 $x_{2i} = \lambda x_{3i}$ 代入：
\begin{itemize}
    \item $\sum y_i x_{2i} = \sum y_i (\lambda x_{3i}) = \lambda \sum y_i x_{3i}$
    \item $\sum x_{2i} x_{3i} = \lambda \sum x_{3i}^2$（刚才算过了）
\end{itemize}

所以分子变成：

\[
(\lambda \sum y_i x_{3i})(\sum x_{3i}^2) - (\sum y_i x_{3i})(\lambda \sum x_{3i}^2)
\]
\[
= \lambda (\sum y_i x_{3i})(\sum x_{3i}^2) - \lambda (\sum y_i x_{3i})(\sum x_{3i}^2)
\]
\[
= 0
\]

\textbf{分子也等于0！}

因此：

\begin{equation}
\hat{\beta}_2 = \frac{0}{0} \tag{4.7}
\end{equation}

这是一个\textbf{不定式}——0/0 可以等于任何数，完全无法确定。

\textbf{再处理 $\hat{\beta}_3$ 的分子}：

分子 $= (\sum y_i x_{3i})(\sum x_{2i}^2) - (\sum y_i x_{2i})(\sum x_{2i} x_{3i})$

把 $x_{2i} = \lambda x_{3i}$ 代入：
\begin{itemize}
    \item $\sum x_{2i}^2 = \lambda^2 \sum x_{3i}^2$
    \item $\sum y_i x_{2i} = \lambda \sum y_i x_{3i}$
    \item $\sum x_{2i} x_{3i} = \lambda \sum x_{3i}^2$
\end{itemize}

所以分子变成：

\[
(\sum y_i x_{3i})(\lambda^2 \sum x_{3i}^2) - (\lambda \sum y_i x_{3i})(\lambda \sum x_{3i}^2)
\]
\[
= \lambda^2 (\sum y_i x_{3i})(\sum x_{3i}^2) - \lambda^2 (\sum y_i x_{3i})(\sum x_{3i}^2)
\]
\[
= 0
\]

\textbf{分子又是0！} 所以：

\begin{equation}
\hat{\beta}_3 = \frac{0}{0} \tag{4.8}
\end{equation}

同样是不定式。

\subsubsection{第四步：直觉理解——为什么估计值会不确定？}

回忆一下 $\hat{\beta}_2$ 的含义：\textbf{在 $X_3$ 保持不变的条件下}，$X_2$ 每变动1单位，$Y$ 平均变动多少。

但现在 $X_2$ 和 $X_3$ 完全绑在一起了（$X_2 = \lambda X_3$）。你根本不可能做到“让 $X_2$ 变而 $X_3$ 不变”——因为 $X_2$ 一变，$X_3$ 必然跟着变！

\textbf{比喻}：假设你要研究“左手力气”和“右手力气”哪个对“举重成绩”影响更大。但数据中所有运动员的左右手力气完全成比例（左手力气 = 2 × 右手力气）。那你根本无法区分：到底是左手还是右手在起作用？因为它们永远一起变，你找不到“右手力气变了但左手力气没变”的情况。两个变量的影响被彻底绑死了，无法分开。

\textbf{从矩阵角度理解}：OLS估计需要计算 $(\mathbf{X}'\mathbf{X})^{-1}$。完全共线性时，$\mathbf{X}$ 矩阵的列线性相关，$\mathbf{X}'\mathbf{X}$ 的行列式为0，逆矩阵不存在。正规方程组的解不唯一——就像方程组 $x + y = 5$ 和 $2x + 2y = 10$ 一样，两个方程其实说的是同一件事，有无穷多组解。

\subsection{后果2：参数估计值的方差无限大}

\textbf{目标}：证明当 $X_2$ 和 $X_3$ 完全共线性时，$\hat{\beta}_2$ 和 $\hat{\beta}_3$ 的方差趋向无穷大。

\subsubsection{第一步：回顾方差公式}

由第三章，两个解释变量模型中 $\hat{\beta}_2$ 和 $\hat{\beta}_3$ 的方差为：

\begin{equation}
\Var(\hat{\beta}_2) = \frac{\sum x_{3i}^2}{(\sum x_{2i}^2)(\sum x_{3i}^2) - (\sum x_{2i} x_{3i})^2} \sigma^2 \tag{4.9}
\end{equation}

\begin{equation}
\Var(\hat{\beta}_3) = \frac{\sum x_{2i}^2}{(\sum x_{2i}^2)(\sum x_{3i}^2) - (\sum x_{2i} x_{3i})^2} \sigma^2 \tag{4.10}
\end{equation}

\textbf{注意}：这两个方差公式的分母与前面估计量公式的分母完全一样！都是 $(\sum x_{2i}^2)(\sum x_{3i}^2) - (\sum x_{2i} x_{3i})^2$。

\subsubsection{第二步：代入完全共线性条件}

已知 $x_{2i} = \lambda x_{3i}$，且我们已算出：
\begin{itemize}
    \item $\sum x_{2i}^2 = \lambda^2 \sum x_{3i}^2$
    \item $\sum x_{2i} x_{3i} = \lambda \sum x_{3i}^2$
\end{itemize}

代入 $\Var(\hat{\beta}_2)$ 的分母：

\[
(\sum x_{2i}^2)(\sum x_{3i}^2) - (\sum x_{2i} x_{3i})^2 = (\lambda^2 \sum x_{3i}^2)(\sum x_{3i}^2) - (\lambda \sum x_{3i}^2)^2
\]
\[
= \lambda^2 (\sum x_{3i}^2)^2 - \lambda^2 (\sum x_{3i}^2)^2 = 0
\]

分母为0！所以：

\[
\Var(\hat{\beta}_2) = \frac{\sum x_{3i}^2}{0} \sigma^2 \to \infty
\]

同理，$\Var(\hat{\beta}_3)$ 的分母也是0：

\[
\Var(\hat{\beta}_3) = \frac{\lambda^2 \sum x_{3i}^2}{0} \sigma^2 \to \infty
\]

\subsubsection{第三步：直觉理解——方差无穷大意味着什么？}

方差衡量的是估计量的“波动程度”。方差无穷大，意味着：每次抽不同的样本，算出来的 $\hat{\beta}_2$ 和 $\hat{\beta}_3$ 可能是天差地别的值——可能很大、可能很小、可能是正的、可能是负的，完全没有稳定性。

\textbf{比喻}：想象一个温度计，每次量同一个温度，读数却在上蹿下跳——有时显示零上50度，有时显示零下30度。这个温度计的“方差”就极大，根本不可信。完全共线性下的OLS估计量就是这样一个“坏温度计”。

\section{第二部分：不完全多重共线性产生的后果}

完全共线性（$X_2 = \lambda X_3$）是一种极端情况，现实中几乎不会出现。\textbf{更常见的情况是不完全共线性}——$X_2$ 和 $X_3$ 之间有很强的近似线性关系，但不是完全的。此时 $(\mathbf{X}'\mathbf{X})^{-1}$ 存在，参数可以估计出来，但会带来一系列问题。

\subsection{后果1：参数估计值的方差与协方差增大}

\textbf{目标}：推导出不完全共线性下，$\hat{\beta}_2$ 和 $\hat{\beta}_3$ 的方差如何随共线性程度增大而增大，最终引入方差膨胀因子（VIF）。

\subsubsection{第一步：设定不完全共线性关系}

不完全共线性意味着 $X_2$ 和 $X_3$ 之间有近似线性关系，但不完全。用离差形式表示为：

\begin{equation}
x_{2i} = \lambda x_{3i} + v_i \tag{4.11}
\end{equation}

其中：
\begin{itemize}
    \item $\lambda \neq 0$：一个非零常数。
    \item $v_i$：随机误差项，表示 $x_{2i}$ 不能被 $\lambda x_{3i}$ 完全解释的那部分。
\end{itemize}

\textbf{关键性质}：$\sum x_{3i} v_i = 0$。这个条件的含义是：$v_i$ 与 $x_{3i}$ 不相关。换句话说，$x_{2i}$ 中能被 $x_{3i}$ 线性解释的部分都归入了 $\lambda x_{3i}$，剩余的部分 $v_i$ 与 $x_{3i}$ 正交（垂直）。

\textbf{暂停理解一下}：对比完全共线性和不完全共线性——
\begin{itemize}
    \item 完全共线性：$x_{2i} = \lambda x_{3i}$，没有 $v_i$ 这一项，$x_2$ 完全由 $x_3$ 决定。
    \item 不完全共线性：$x_{2i} = \lambda x_{3i} + v_i$，多了一个 $v_i$。$v_i$ 就是 $x_2$ 中“不能被 $x_3$ 解释”的那部分独立信息。$v_i$ 越小，说明 $x_2$ 能被 $x_3$ 解释的部分越多，共线性越严重。
\end{itemize}

\subsubsection{第二步：将不完全共线性代入估计公式}

将式(4.11)代入式(4.6)，可以验证 $\hat{\beta}_3$ 仍然是可计算的（因为分母不再为0了）。具体来说，把 $x_{2i} = \lambda x_{3i} + v_i$ 代入后，利用 $\sum x_{3i} v_i = 0$ 化简，得到：

\begin{equation}
\hat{\beta}_3 = \frac{(\sum y_i x_{3i})(\lambda^2 \sum x_{3i}^2 + \sum v_i^2) - (\lambda \sum y_i x_{3i} + \sum v_i y_i)(\lambda \sum x_{3i}^2)}{(\lambda^2 \sum x_{3i}^2 + \sum v_i^2)(\sum x_{3i}^2) - \lambda^2 (\sum x_{3i}^2)^2} \tag{4.12}
\end{equation}

\textbf{关键观察}：这个分母不再为0了！因为分子多了 $\sum v_i^2$ 项，使得分母中出现了 $(\lambda^2 \sum x_{3i}^2 + \sum v_i^2)(\sum x_{3i}^2) - \lambda^2 (\sum x_{3i}^2)^2 = (\sum v_i^2)(\sum x_{3i}^2) > 0$。所以 $\hat{\beta}_3$ 是可以算出来的。

\textbf{但是}：如果 $X_2$ 与 $X_3$ 共线程度越高，$v_i$ 就越小（因为 $x_2$ 越来越接近 $\lambda x_3$，$v_i$ 是剩下的“残余”部分）。当 $v_i \to 0$ 时，分母趋近于0，$\hat{\beta}_3$ 趋于不确定——也就是向完全共线性的情况靠拢。

\subsubsection{第三步：引入相关系数 $r_{23}$}

为了量化“共线性有多严重”，我们使用 $X_2$ 和 $X_3$ 之间的样本相关系数的平方。用离差形式表示：

\begin{equation}
r_{23}^2 = \frac{(\sum x_{2i} x_{3i})^2}{(\sum x_{2i}^2)(\sum x_{3i}^2)} \tag{4.13}
\end{equation}

\textbf{暂停理解一下}：$r_{23}^2$ 衡量的是 $X_2$ 和 $X_3$ 之间线性关系的强弱。
\begin{itemize}
    \item $r_{23}^2 = 0$：$X_2$ 和 $X_3$ 毫无线性关系，没有共线性。
    \item $r_{23}^2 = 1$：$X_2$ 和 $X_3$ 完全线性相关，完全共线性。
    \item $r_{23}^2$ 越接近1，共线性越严重。
\end{itemize}

\subsubsection{第四步：推导 $\Var(\hat{\beta}_2)$ 的简化表达式}

现在，我们要把式(4.9)中方差的分母用 $r_{23}^2$ 来改写，使得共线性的影响清晰可见。

原始公式：

\begin{equation}
\Var(\hat{\beta}_2) = \frac{\sum x_{3i}^2}{(\sum x_{2i}^2)(\sum x_{3i}^2) - (\sum x_{2i} x_{3i})^2} \sigma^2 \tag{4.9}
\end{equation}

\textbf{化简步骤}（极其详细，无跳步）：

步骤1：观察分母。分母 $= (\sum x_{2i}^2)(\sum x_{3i}^2) - (\sum x_{2i} x_{3i})^2$。

步骤2：提取公因子 $(\sum x_{2i}^2)(\sum x_{3i}^2)$。也就是把分母写成：

\[
(\sum x_{2i}^2)(\sum x_{3i}^2)\left[1 - \frac{(\sum x_{2i} x_{3i})^2}{(\sum x_{2i}^2)(\sum x_{3i}^2)}\right]
\]

\textbf{暂停理解一下}：这一步是怎么做的？其实就是把 $a \cdot b - c^2$ 写成 $a \cdot b \cdot (1 - c^2/(ab))$。你可以自己验证：$ab(1 - c^2/(ab)) = ab - ab \cdot c^2/(ab) = ab - c^2$。没错吧！

步骤3：认出方括号里就是 $1 - r_{23}^2$！因为根据式(4.13)：

\[
r_{23}^2 = \frac{(\sum x_{2i} x_{3i})^2}{(\sum x_{2i}^2)(\sum x_{3i}^2)}
\]

所以分母变成了：

\[
(\sum x_{2i}^2)(\sum x_{3i}^2)(1 - r_{23}^2)
\]

步骤4：代回方差公式。现在整个方差变为：

\[
\Var(\hat{\beta}_2) = \frac{\sum x_{3i}^2}{(\sum x_{2i}^2)(\sum x_{3i}^2)(1 - r_{23}^2)} \sigma^2
\]

步骤5：约分。分子有 $\sum x_{3i}^2$，分母也有 $\sum x_{3i}^2$，约掉：

\begin{equation}
\Var(\hat{\beta}_2) = \frac{\sigma^2}{\sum x_{2i}^2 (1 - r_{23}^2)} \tag{4.14}
\end{equation}

\subsubsection{第五步：同理推导 $\Var(\hat{\beta}_3)$}

用完全对称的方式（分子换为 $\sum x_{2i}^2$，约分时约掉 $\sum x_{2i}^2$），得到：

\begin{equation}
\Var(\hat{\beta}_3) = \frac{\sigma^2}{\sum x_{3i}^2 (1 - r_{23}^2)} \tag{4.15}
\end{equation}

\subsubsection{第六步：推导协方差}

类似地，可以推导出：

\begin{equation}
\Cov(\hat{\beta}_2, \hat{\beta}_3) = \frac{-r_{23} \sigma^2}{(1 - r_{23}^2)\sqrt{(\sum x_{2i}^2)(\sum x_{3i}^2)}} \tag{4.16}
\end{equation}

\textbf{暂停理解一下}：注意协方差前面的负号和 $r_{23}$。当 $r_{23} > 0$（$X_2$ 和 $X_3$ 正相关）时，协方差为负。直觉上：如果 $X_2$ 和 $X_3$ 同向变动，那么你高估 $\beta_2$ 的同时往往低估 $\beta_3$，两个估计量呈反向变动，协方差为负。

\subsubsection{第七步：关键观察——$r_{23}^2$ 越大，方差越大}

现在看式(4.14)和(4.15)，分母中都有 $(1 - r_{23}^2)$ 这一项。

\begin{itemize}
    \item 如果 $r_{23}^2 = 0$（无共线性）：分母中的 $(1-0)=1$，方差就是 $\sigma^2/\sum x_{2i}^2$，这是“正常的”方差。
    \item 如果 $r_{23}^2 = 0.5$：$(1-0.5)=0.5$，方差变为无共线性时的 $1/0.5 = 2$ 倍。
    \item 如果 $r_{23}^2 = 0.9$：$(1-0.9)=0.1$，方差变为 $1/0.1 = 10$ 倍！
    \item 如果 $r_{23}^2 = 0.99$：$(1-0.99)=0.01$，方差变为 $1/0.01 = 100$ 倍！！
    \item 如果 $r_{23}^2 \to 1$：$(1-r_{23}^2) \to 0$，方差 $\to \infty$——回到了完全共线性的结果。
\end{itemize}

\textbf{比喻}：$(1-r_{23}^2)$ 就像是水管的水压。$r_{23}^2$ 越大，水管越窄（水压越低），方差（用水困难程度）就越大。当水管完全堵死（$r_{23}^2 = 1$），水完全流不过去（方差无穷大）。

\subsubsection{第八步：引入方差膨胀因子（VIF）}

为了更直观地衡量共线性对方差的影响，我们定义\textbf{方差膨胀因子}（Variance Inflation Factor，VIF）：

\begin{equation}
\VIF = \frac{1}{1 - r_{23}^2} \tag{4.17}
\end{equation}

\textbf{为什么叫“膨胀因子”？} 因为它衡量的是：由于共线性的存在，方差“膨胀”了多少倍。

\begin{itemize}
    \item 当 $r_{23}^2 = 0$（无共线性）：VIF $= 1/(1-0) = 1$。方差没有膨胀。
    \item 当 $r_{23}^2 = 0.9$：VIF $= 1/0.1 = 10$。方差膨胀了10倍。
    \item 当 $r_{23}^2 = 0.99$：VIF $= 1/0.01 = 100$。方差膨胀了100倍！
    \item 当 $r_{23}^2 \to 1$：VIF $\to \infty$。方差膨胀到无穷大。
\end{itemize}

用VIF重写方差公式：

\begin{equation}
\Var(\hat{\beta}_2) = \frac{\sigma^2}{\sum x_{2i}^2} \times \VIF \tag{4.18}
\end{equation}

\begin{equation}
\Var(\hat{\beta}_3) = \frac{\sigma^2}{\sum x_{3i}^2} \times \VIF \tag{4.19}
\end{equation}

\textbf{极其重要的结论}：$\hat{\beta}_2$ 和 $\hat{\beta}_3$ 的方差与VIF成正比。VIF越大，方差越大，估计越不精确。

\textbf{暂停理解一下}：$\sigma^2/\sum x_{2i}^2$ 是什么？它是在没有共线性（$r_{23}^2=0$）时 $\hat{\beta}_2$ 的“基准方差”。而VIF就是这个基准方差被“放大”的倍数。所以VIF就是“膨胀倍数”。

\subsection{后果2：置信区间扩大}

\textbf{目标}：说明共线性如何使得参数的置信区间变宽，降低估计精度。

\subsubsection{第一步：回顾置信区间的构造}

$\beta_3$ 的95\%置信区间为：

\[
\hat{\beta}_3 \pm 1.96 \sqrt{\Var(\hat{\beta}_3)}
\]

其中1.96是标准正态分布在95\%置信度下的双侧分位数。

将式(4.15)代入：

\[
\hat{\beta}_3 \pm 1.96 \sqrt{\frac{\sigma^2}{\sum x_{3i}^2 (1 - r_{23}^2)}} = \hat{\beta}_3 \pm 1.96 \sqrt{\frac{\sigma^2}{\sum x_{3i}^2}} \times \frac{1}{\sqrt{1 - r_{23}^2}}
\]

注意 $\frac{1}{\sqrt{1-r_{23}^2}} = \sqrt{\VIF}$。所以置信区间的半宽度（即 $\pm$ 后面的部分）与 $\sqrt{\VIF}$ 成正比！

\textbf{暂停理解一下}：置信区间越窄，我们对参数真实值的定位越精确；越宽，越不确定。共线性通过增大VIF，使得置信区间变宽，我们就像在更大范围内“瞎猜”参数在哪里。

\subsubsection{第二步：表4.1的数值演示}

书中的表4.1给出了不同 $r_{23}$ 下 $\beta_3$ 的95\%置信区间：

\begin{table}[htbp]
\centering
\caption{表4.1 不同$r_{23}$下$\beta_3$的95\%置信区间}
\small
\begin{tabular}{ccccc}
\toprule
\textbf{$r_{23}$} & \textbf{$r_{23}^2$} & \textbf{VIF} & \textbf{$\beta_3$的95\%置信区间} & \textbf{相对于无共线性的宽度倍数} \\
\midrule
0.00 & 0.00 & 1 & $\hat{\beta}_3 \pm 1.96\sqrt{\sigma^2/\sum x_{3i}^2}$ & 1倍 \\
0.50 & 0.25 & 1.33 & $\hat{\beta}_3 \pm 1.96\sqrt{1.33\sigma^2/\sum x_{3i}^2}$ & $\sqrt{1.33} \approx 1.15$倍 \\
0.99 & 0.9801 & 50.25 & $\hat{\beta}_3 \pm 1.96\sqrt{50\sigma^2/\sum x_{3i}^2}$ & $\sqrt{50} \approx 7.07$倍 \\
0.999 & 0.998001 & 500.25 & $\hat{\beta}_3 \pm 1.96\sqrt{500\sigma^2/\sum x_{3i}^2}$ & $\sqrt{500} \approx 22.4$倍 \\
\bottomrule
\end{tabular}
\end{table}

\textbf{逐行解释}：

\begin{itemize}
    \item \textbf{$r_{23} = 0$}：没有共线性，VIF = 1，这是“正常的”置信区间宽度。
    \item \textbf{$r_{23} = 0.50$}：VIF $= 1/(1-0.25) = 1/0.75 \approx 1.33$。方差只膨胀了1.33倍，置信区间宽度约为原来的1.15倍——影响不大。
    \item \textbf{$r_{23} = 0.99$}：VIF $= 1/(1-0.9801) = 1/0.0199 \approx 50.25$。方差膨胀了50倍！置信区间宽度约为原来的 $\sqrt{50} \approx 7$ 倍——影响巨大！
    \item \textbf{$r_{23} = 0.999$}：VIF $\approx 500$。方差膨胀了500倍！置信区间宽度约为原来的 $\sqrt{500} \approx 22$ 倍——完全不可用了！
\end{itemize}

\textbf{比喻}：你在玩飞镖。无共线性时，你大概能扎在靶心附近（小置信区间）。$r_{23}=0.99$ 时，就像蒙着眼睛转了7圈再扔飞镖——飞镖可能扎到靶上任何地方（大置信区间）。$r_{23}=0.999$ 时，你不仅蒙着眼睛转了22圈，还被人扛到了另一个房间——飞镖可能扎到墙上、天花板上，完全不可控。

\subsection{后果3：假设检验容易做出错误的判断}

\textbf{目标}：说明共线性如何导致t检验失效——本应显著的变量被误判为不显著。

\subsubsection{第一步：t检验的统计量}

对回归系数的原假设 $H_0: \beta_3 = 0$ 进行检验时，使用t统计量：

\[
t = \frac{\hat{\beta}_3}{\sqrt{\Var(\hat{\beta}_3)}}
\]

\textbf{暂停理解一下}：t值的含义是“估计值是标准误的几倍”。t值越大，说明估计值相对于其波动幅度而言“很大”，越有理由认为 $\beta_3$ 不是0（即 $X_3$ 对 $Y$ 有显著影响）。

\subsubsection{第二步：共线性如何影响t值}

由式(4.15)：

\[
\Var(\hat{\beta}_3) = \frac{\sigma^2}{\sum x_{3i}^2 (1 - r_{23}^2)}
\]

当 $r_{23}^2$ 增大时，$(1-r_{23}^2)$ 减小，$\Var(\hat{\beta}_3)$ 增大，$\sqrt{\Var(\hat{\beta}_3)}$ 也增大。

现在看t值：

\[
t = \frac{\hat{\beta}_3}{\sqrt{\Var(\hat{\beta}_3)}}
\]

分母变大，t值变小。

\subsubsection{第三步：t值变小会导致什么？}

t检验的判断规则是：如果 $|t|$ 大于临界值（比如1.96），就拒绝 $H_0: \beta_3 = 0$，认为 $X_3$ 对 $Y$ 有显著影响。

但共线性使得t值变小。原本 $|t|$ 应该大于1.96（$X_3$ 确实有影响），现在却因为共线性导致t值缩小到1.96以下，于是我们\textbf{错误地}接受了“$\beta_3 = 0$”的原假设。

\textbf{比喻}：你在法庭上想证明某人有罪。证据（$\hat{\beta}_3$）确实指向他有罪，但因为你对证据的“不确定性”太大（方差太大、标准误太大），法官说“证据不够确凿”，放了他。实际上他是有罪的，只是共线性干扰了你的判断。

\textbf{双重打击}：
\begin{enumerate}
    \item 置信区间扩大 $\to$ 更容易接受本应拒绝的假设（第一类错误的风险变化）。
    \item t值变小 $\to$ 更容易错误地接受“$\beta = 0$”的原假设（第二类错误增大）。
\end{enumerate}

\subsection{后果4：$R^2$ 高但 $t$ 检验不显著——最令人困惑的现象}

\textbf{目标}：解释为什么会出现“整体方程很显著，但单个变量都不显著”的矛盾现象。

\subsubsection{第一步：矛盾现象描述}

在存在严重多重共线性时，常常出现以下情况：

\begin{itemize}
    \item \textbf{$R^2$ 很高}：模型的拟合优度很好，解释变量“集体”对 $Y$ 的解释力很强。
    \item \textbf{$F$ 检验显著}：联合检验拒绝“所有系数都为0”的原假设，说明解释变量整体对 $Y$ 有显著影响。
    \item \textbf{但个别 $t$ 检验不显著}：对 $\beta_2$ 和 $\beta_3$ 分别做 $t$ 检验，却都接受“$\beta = 0$”的原假设。
    \item \textbf{甚至系数符号反常}：比如经济理论上 $\beta_2$ 应该为正，但估计出来却是负的。
\end{itemize}

\subsubsection{第二步：为什么会出现这种矛盾？}

\textbf{直觉解释}：

\begin{itemize}
    \item $R^2$ 高、$F$ 检验显著，说明 $X_2$ 和 $X_3$ \textbf{合在一起}确实能很好地解释 $Y$。这没问题——因为它们的信息总量是够的。
    \item 但 $t$ 检验不显著，是因为我们无法\textbf{单独}区分 $X_2$ 和 $X_3$ 各自的贡献。它们的信息高度重叠，就像两个人紧挨着站在一起，你能看到“有人在那”，但分不清谁是谁。
    \item 系数符号反常，是因为共线性导致估计值极不稳定。就像温度计坏了一样，它可能给你一个完全相反的读数。
\end{itemize}

\textbf{比喻}：两个歌手合唱一首歌，合唱效果非常好（$R^2$高），但你听不出每个人唱了多大声——因为他们的声音几乎完全同步（$t$检验不显著）。甚至有时候因为回声干扰，你误以为其中一个人在唱反调（系数符号反常）。

\subsubsection{第三步：教材中的实例}

教材提到了一个典型例子：“工业增加值增长会减少财政收入吗？”

正常情况下，工业增加值增长，财政收入应该增加（$\beta$ 应该为正）。但由于工业增加值与其他解释变量之间存在严重共线性，回归结果中工业增加值的系数竟然为负——这在经济学上完全说不通。这正是多重共线性导致的“系数符号反常”现象。

\textbf{遇到这种情况时}：如果发现回归方程整体很显著（$R^2$高、$F$检验通过），但个别系数 $t$ 检验不显著、甚至符号与经济理论矛盾，\textbf{首先应该怀疑的就是多重共线性}。

\section{总结：多重共线性的全部后果一览}

严重的多重共线性常常导致以下情形出现：

\textbf{1. 估计值不稳定}：用OLS得到的回归参数估计值很不稳定，换一个样本可能得到完全不同的结果。

\textbf{2. 方差与协方差加速增长}：回归系数的方差与协方差随着多重共线性强度的增加而加速增长（不是线性增长，而是通过VIF的机制加速膨胀），对参数难以做出精确的估计。

\textbf{3. 整体显著但个别不显著}：造成回归方程高度显著（$R^2$高、$F$检验通过）的情况下，有些回归系数却通不过显著性检验（$t$检验不通过）。

\textbf{4. 系数符号反常}：甚至可能出现回归系数的正负号得不到合理的经济解释。

\textbf{但是——有一个重要的例外}：如果研究目的\textbf{仅在于预测 $Y$}，而且各个解释变量 $X$ 之间的多重共线性关系的性质在未来将继续保持，这时虽然无法精确估计个别的回归系数，但可以估计这些系数的某些线性组合，因此多重共线性可能并不是严重问题。

\textbf{暂停理解一下这最后一条}：为什么预测时共线性不是大问题？因为预测只关心 $\hat{Y} = \hat{\beta}_1 + \hat{\beta}_2 X_2 + \hat{\beta}_3 X_3$ 这个整体，而不关心 $\hat{\beta}_2$ 和 $\hat{\beta}_3$ 各自是多少。虽然 $\hat{\beta}_2$ 和 $\hat{\beta}_3$ 各自不稳定，但它们“此消彼长”（协方差为负），加在一起的效果可能是稳定的——只要 $X_2$ 和 $X_3$ 之间的关系在未来保持不变。

\textbf{比喻}：你要预测明天的天气，用“温度计A”和“温度计B”两个指标。虽然这两个指标高度相关，你分不清各自对天气预测的独立贡献，但只要它们的相关关系不变，你用它们一起预测天气的效果仍然很好——只是你无法说“温度计A贡献了60\%，温度计B贡献了40\%”。

\vspace{1em}
\noindent\textbf{整篇文章的逻辑链条}：

\begin{enumerate}
    \item \textbf{完全共线性}：OLS估计量的分母为0 → 估计值不确定（0/0）；方差公式的分母也为0 → 方差无穷大。核心原因：两个变量提供的是完全相同的信息，无法区分各自的贡献。
    \item \textbf{不完全共线性}：估计量可以计算，但共线性越强（$r_{23}^2$ 越大），VIF越大 → 方差越大 → 置信区间越宽 → t值越小 → 假设检验越容易犯错 → 可能出现“整体显著但个别不显著”甚至“系数符号反常”的矛盾现象。
    \item \textbf{例外情况}：如果目的只是预测，且共线性结构未来保持，则多重共线性不是大问题。
\end{enumerate}

希望这样的讲解能够让你完全明白每一步是在做什么、为什么这样做。

\end{document}